\documentclass[11pt,a4paper]{article}
\usepackage[utf8]{vietnam}
\usepackage{amsmath,amssymb}
\usepackage{tabularx}
\usepackage{tikz}
\usepackage{xcolor}
\usepackage{fancyhdr}
\usepackage[top=1.8cm, bottom=2.0cm, left=1.8cm, right=1.8cm]{geometry}

\setlength{\parskip}{3pt}
\setlength{\parindent}{0pt}

\pagestyle{fancy}
\fancyhf{}
\renewcommand{\headrulewidth}{0.4pt}
\lhead{\footnotesize \textbf{CHÙA HƯƠNG SƠN -- ÔN THI VÀO 10}}
\rhead{\footnotesize \textbf{ĐÁP ÁN \& LỜI GIẢI CHI TIẾT (ĐỀ SỐ 08)}}
\cfoot{\thepage}

\begin{document}

\begin{center}
\textbf{\large HƯỚNG DẪN CHẤM VÀ ĐÁP ÁN ĐỀ THI THỬ SỐ 08}\\
\textbf{MÔN: TOÁN 9 -- HỌC KỲ I -- NĂM HỌC 2026 -- 2027}\\
\textit{(Hướng dẫn chấm gồm 02 trang)}
\end{center}

\vspace{1mm}
\noindent
\rule{\linewidth}{0.8pt}

\vspace{2mm}
\noindent
\textbf{\large A. ĐÁP ÁN TRẮC NGHIỆM (7,0 ĐIỂM)}

\vspace{2mm}
\noindent
\textbf{1. Phần 1: Trắc nghiệm nhiều lựa chọn (3,0 điểm -- Mỗi câu đúng 0,25 điểm)}
\begin{center}
\begin{tabular}{|c|c|c|c|c|c|c|c|c|c|c|c|}
\hline
\textbf{1} & \textbf{2} & \textbf{3} & \textbf{4} & \textbf{5} & \textbf{6} & \textbf{7} & \textbf{8} & \textbf{9} & \textbf{10} & \textbf{11} & \textbf{12} \\
\hline
\textbf{A} & \textbf{A} & \textbf{A} & \textbf{A} & \textbf{A} & \textbf{A} & \textbf{A} & \textbf{A} & \textbf{A} & \textbf{A} & \textbf{A} & \textbf{A} \\
\hline
\end{tabular}
\end{center}

\vspace{2mm}
\noindent
\textbf{2. Phần 2: Trắc nghiệm Đúng -- Sai (2,0 điểm)}
\begin{itemize}
  \item \textbf{Câu 13 (1,0 điểm):} a) Đúng \quad b) Đúng \quad c) Đúng \quad d) Sai.
  \item \textbf{Câu 14 (1,0 điểm):} a) Đúng \quad b) Đúng \quad c) Sai \quad d) Đúng.
\end{itemize}
\textit{Quy tắc chấm: Đúng 1 ý được 0,1đ; đúng 2 ý được 0,25đ; đúng 3 ý được 0,5đ; đúng 4 ý được 1,0đ.}

\vspace{2mm}
\noindent
\textbf{3. Phần 3: Trắc nghiệm trả lời ngắn (2,0 điểm -- Mỗi câu đúng 0,5 điểm)}
\begin{itemize}
  \item \textbf{Câu 15:} $x < 3$.
  \item \textbf{Câu 16:} $3$.
  \item \textbf{Câu 17:} $29{,}5\text{ m}$ (hoặc $29{,}5$).
  \item \textbf{Câu 18:} $18{,}84\text{ m}^2$ (hoặc $18{,}84$).
\end{itemize}

\vspace{3mm}
\noindent
\textbf{\large B. LỜI GIẢI CHI TIẾT VÀ BAREM ĐIỂM TỰ LUẬN (3,0 ĐIỂM)}

\vspace{2mm}
\noindent
\textbf{Câu 19 (1,0 điểm).}
\begin{itemize}
  \item[a)] \textit{(0,5 điểm)} Hệ phương trình: $\begin{cases} 4x - 3y = 7 \quad (1) \\ 3x + 2y = 1 \quad (2) \end{cases}$.\\
  Nhân phương trình $(1)$ với $2$ và nhân phương trình $(2)$ với $3$, ta được:
  $\begin{cases} 8x - 6y = 14 \\ 9x + 6y = 3 \end{cases}$ \dotfill \textbf{(0,25đ)}\\
  Cộng hai phương trình vế theo vế: $17x = 17 \iff x = 1$.\\
  Thay $x = 1$ vào phương trình $(2)$: $3(1) + 2y = 1 \iff 2y = -2 \iff y = -1$.\\
  Vậy hệ phương trình có nghiệm duy nhất là $(x; y) = (1; -1)$. \dotfill \textbf{(0,25đ)}

  \item[b)] \textit{(0,5 điểm)} Với $x \ge 0, x \ne 9$, ta có:
  \[
  \dfrac{\sqrt{x}}{\sqrt{x}-3} - \dfrac{3}{\sqrt{x}+3} = \dfrac{\sqrt{x}(\sqrt{x}+3) - 3(\sqrt{x}-3)}{(\sqrt{x}-3)(\sqrt{x}+3)} = \dfrac{x + 3\sqrt{x} - 3\sqrt{x} + 9}{x - 9} = \dfrac{x+9}{x-9} \quad \text{\textbf{(0,25đ)}}
  \]
  Do đó:
  \[
  B = \dfrac{x+9}{x-9} : \dfrac{x+9}{x-9} = \dfrac{x+9}{x-9} \cdot \dfrac{x-9}{x+9} = 1 \quad \text{\textbf{(0,25đ)}}
  \]
\end{itemize}

\vspace{2mm}
\noindent
\textbf{Câu 20 (1,5 điểm).}
\begin{itemize}
  \item[a)] \textit{(0,75 điểm)}
  Theo tính chất của hai tiếp tuyến cắt nhau:
  $CA = CM$ và $DB = DM$. Suy ra $CD = CM + DM = CA + BD$ (đpcm). \dotfill \textbf{(0,25đ)}\\
  Mặt khác, $OC$ là tia phân giác của $\widehat{AOM}$ nên $\widehat{COM} = \dfrac{1}{2}\widehat{AOM}$.\\
  $OD$ là tia phân giác của $\widehat{BOM}$ nên $\widehat{DOM} = \dfrac{1}{2}\widehat{BOM}$. \dotfill \textbf{(0,25đ)}\\
  Do hai góc $\widehat{AOM}$ và $\widehat{BOM}$ kề bù nên:
  $\widehat{COD} = \widehat{COM} + \widehat{DOM} = \dfrac{1}{2}(\widehat{AOM} + \widehat{BOM}) = \dfrac{1}{2} \cdot 180^\circ = 90^\circ$ (đpcm). \dotfill \textbf{(0,25đ)}

  \item[b)] \textit{(0,75 điểm)}
  Xét tam giác $COD$ vuông tại $O$, có đường cao $OM \perp CD$ (vì $CD$ là tiếp tuyến tại $M$).\\
  Áp dụng hệ thức lượng: $CM \cdot MD = OM^2 = R^2$.\\
  Mà $CM = AC$ và $MD = BD$, suy ra $AC \cdot BD = R^2$ (đpcm). \dotfill \textbf{(0,5đ)}\\
  Gọi $I$ là trung điểm của $CD$. Khi đó $I$ là tâm của đường tròn đường kính $CD$, bán kính $r = \dfrac{CD}{2}$.\\
  Tứ giác $ABDC$ là hình thang vuông ($AC \parallel BD \perp AB$) có $O$ là trung điểm $AB$ và $I$ là trung điểm $CD$, nên $OI$ là đường trung bình của hình thang $ABDC$.\\
  Suy ra $OI \parallel AC \implies OI \perp AB$ tại $O$, và $OI = \dfrac{AC + BD}{2} = \dfrac{CD}{2} = r$.\\
  Vậy đường tròn $(I; r)$ đường kính $CD$ tiếp xúc với đường thẳng $AB$ tại $O$. \dotfill \textbf{(0,25đ)}
\end{itemize}

\vspace{2mm}
\noindent
\textbf{Câu 21 (0,5 điểm).}
\begin{itemize}
  \item[a)] \textit{(0,25 điểm)}
  Thể tích bể hình trụ không nắp là $V = \pi R^2 h = 27\pi \implies h = \dfrac{27}{R^2}$.\\
  Diện tích tôn cần dùng (diện tích toàn phần không nắp) là:
  \[
  S = \pi R^2 + 2\pi R h = \pi R^2 + 2\pi R \cdot \dfrac{27}{R^2} = \pi \left( R^2 + \dfrac{54}{R} \right) = \pi \left( R^2 + \dfrac{27}{R} + \dfrac{27}{R} \right).
  \]
  Áp dụng bất đẳng thức Cô-si cho 3 số dương $R^2, \dfrac{27}{R}, \dfrac{27}{R}$:
  \[
  R^2 + \dfrac{27}{R} + \dfrac{27}{R} \ge 3\sqrt[3]{R^2 \cdot \dfrac{27}{R} \cdot \dfrac{27}{R}} = 3\sqrt[3]{729} = 3 \cdot 9 = 27.
  \]
  Dấu ``='' xảy ra khi $R^2 = \dfrac{27}{R} \iff R^3 = 27 \iff R = 3\text{ (m)}$. Khi đó $h = \dfrac{27}{3^2} = 3\text{ (m)}$.\\
  Vậy để tốn ít tôn nhất thì $R = 3\text{ m}$ và $h = 3\text{ m}$. \dotfill \textbf{(0,25đ)}

  \item[b)] \textit{(0,25 điểm)}
  Với $x, y > 0$, ta có:
  \[
  P = \dfrac{1}{x^2 + y^2} + \dfrac{1}{xy} = \left( \dfrac{1}{x^2 + y^2} + \dfrac{1}{2xy} \right) + \dfrac{1}{2xy}.
  \]
  Áp dụng bất đẳng thức $\dfrac{1}{a} + \dfrac{1}{b} \ge \dfrac{4}{a+b}$ cho hai số dương $x^2+y^2$ và $2xy$:
  \[
  \dfrac{1}{x^2 + y^2} + \dfrac{1}{2xy} \ge \dfrac{4}{x^2 + y^2 + 2xy} = \dfrac{4}{(x+y)^2} = \dfrac{4}{2^2} = 1.
  \]
  Mặt khác, $xy \le \dfrac{(x+y)^2}{4} = \dfrac{4}{4} = 1 \implies \dfrac{1}{2xy} \ge \dfrac{1}{2 \cdot 1} = \dfrac{1}{2}$.\\
  Suy ra: $P \ge 1 + \dfrac{1}{2} = \dfrac{3}{2}$.\\
  Dấu đẳng thức xảy ra khi $\begin{cases} x^2 + y^2 = 2xy \\ x = y \\ x + y = 2 \end{cases} \iff x = y = 1$.\\
  Vậy giá trị nhỏ nhất của $P$ bằng $\dfrac{3}{2}$, đạt được khi $x = y = 1$. \dotfill \textbf{(0,25đ)}
\end{itemize}

\end{document}
